\documentclass[10pt,a4paper]{amsart}
\usepackage[margin=19mm]{geometry}
\usepackage{amsmath,amssymb,mathtools}
\usepackage[scr=rsfs,cal=euler]{mathalfa}
\usepackage{xfrac}
\usepackage[unicode,hidelinks,bookmarks=false]{hyperref}
\newcommand{\ie}{i.e.\ }
\newcommand{\cf}{cf.\ }
\newcommand{\resp}{respectively\ }
\newcommand{\Subsection}{\subsection}
\providecommand{\nobreakdash}{\nobreak\hbox{-}}
\setlength{\emergencystretch}{2em}
\multlinegap0pt
\usepackage{amssymb}
\usepackage{tikz}
\usepackage{dsfont}
\usepackage{cancel}
\usepackage{xcolor}
\usepackage{mathtools}
\newtheorem{proposition}{Proposition}
\newtheorem{theorem}{Theorem}
\title{On Local Finiteness of Modal K4 Algebras}
\author{Gabriel Agnew}
\date{Source: arXiv:2606.02941v1 (CC BY 4.0)}
\begin{document}
\maketitle

\noindent\emph{Source and licence.} On Local Finiteness of Modal K4 Algebras. Version arXiv:2606.02941v1 (CC BY 4.0), distributed by arXiv under the Creative Commons Attribution 4.0 International licence, http://creativecommons.org/licenses/by/4.0/, as recorded in the arXiv licence metadata for this exact version and shown on the abstract page https://arxiv.org/abs/2606.02941v1. Full licence notice: \texttt{licenses/arxiv-2606.02941-CC-BY-4.0.txt}.\par\smallskip

\noindent\textit{Editorial scope.} Counted unit: the article's theorem thm:no\_tower\_implies\_locally\_finite (source0.tex lines 218--221). The article's complete original proof of it is delivered with it (source0.tex lines 223--329, one \texttt{proof} environment of 107 lines). No mathematics is written, repaired or re-derived: the statement and the proof are the article's own lines, in the article's own notation, and no retained line is substituted or re-flowed.  Retained uncounted as the source-local context the proof needs: lines 161--164; lines 176--202.  Nothing was omitted from any retained span, so the package carries no omission record.  No retained line contains a citation, so the wrapper carries no bibliography and no bibliography entry.  No retained line contains a figure environment, an image-inclusion command or drawing code, so no image is delivered and the package declares no asset dependency.  Editorial formatting only, in addition: an AMS \texttt{amsart} preamble that restates the article's own declarations and macros used by the retained lines, namely the environments \texttt{proposition}, \texttt{theorem} on separate counters, together with the standard packages those lines need.  The retained environments print in this document as Proposition 1, Theorem 1 in order of appearance, whereas the article prints the same items with the numbering of its own full document.  The complete native source of the article as distributed in the arXiv e-print for this exact version is bundled unchanged as \texttt{sources/201944/source0.tex}.
\begin{proposition} \label{prop:splitting}
    Let $(X,R,V)$ be a general frame, and $\mathcal{A},\mathcal{B} \subseteq V$ be two finite nonempty families of subsets of $X$ which partition $\bigcup \mathcal{A}$ and $\bigcup \mathcal{B}$, respectively.
    Then the splitting of $\mathcal{A}$ among $\mathcal{B}$ is a finite partition of $\bigcup \mathcal{A}$ which refines $\mathcal{A}$ and is contained in $V$.
\end{proposition}

\begin{proposition}\label{prop:exists_tuned_wrt}
    Let $(X,R,V)$ be a general frame 
    with $\mathcal{A},\mathcal{C} \subseteq V$ finite nonempty families of subsets of $X$ which partition $\bigcup\mathcal{A}$ and $\bigcup \mathcal{C}$ respectively.
    Then there is a finite refinement $\mathcal{A}'\subseteq V$ of $\mathcal{A}$ which is tuned with respect to $\mathcal{C}$.
    

\begin{proof}
    If $\bigcup\mathcal{A} \cap \Diamond\bigcup \mathcal{C}=\emptyset$, we are done. So, suppose not.
    Then for each $x \in \bigcup \mathcal{A}$, there is a collection $\mathcal{W}\subseteq\mathcal{C}$ such that $x \in \Diamond W$ for each $W\in \mathcal{W}$ and $x\notin \Diamond U$ for any $U \in \mathcal{C-W}$.
    Ignoring instances of $\emptyset$, split $\mathcal{A}$ over the (nonempty) family of disjoint subsets
    \[
    \left\{\bigcap_{W \in \mathcal{W}} \Diamond W \cap 
    \bigcap_{U\in (\mathcal{C-W})}\Box (X-U) 
    \mid \emptyset \neq \mathcal{W} \subseteq \mathcal{C}\right\} 
    \]
    and call the result of this splitting $\mathcal{A}'$. Since the above family is a finite, nonempty collection of members of $V$, from proposition \ref{prop:splitting} it follows that $\mathcal{A}' \subseteq V$ and is finite. 
    Now, let $\mathcal{B}$ be any refinement of $\mathcal{A}'$, and let $B \in \mathcal{B}$, $C \in \mathcal{C}$.
    Suppose $bRc$ for some $b \in B$, $c \in C$. Then 
    $b \in \Diamond C$, so there is a nonempty $\mathcal{W} \subseteq \mathcal{C}$ such that $C \in \mathcal{W}$ and
    \[
    b \in \bigcap_{W \in \mathcal{W}} \Diamond W \cap \bigcap_{U\in (\mathcal{C-W})}\Box (X-U).
    \]
    But any $d \in B$ will also be a member of this set. Hence $d \in \Diamond C$ as well, and it follows that $\mathcal{B}$ is tuned with respect to $\mathcal{C}$.
    
\end{proof}
    
\end{proposition}

\begin{theorem} \label{thm:no_tower_implies_locally_finite}
    Let $(V,\Diamond)$ be a $K4$ modal algebra.
    Then if $(V,\Diamond)$ does not contain an infinite descending $\Diamond$-tower, it is locally finite.
\end{theorem}

\begin{proof}
    Let $(X,R,V)$ be any general frame whose dual modal algebra is $(V,\Diamond)$, and assume $(V,\Diamond)$ has no infinite descending $\Diamond$-towers. 
    Let $P_1 \subseteq V$ be a finite partition of $X$.
    We construct a refinement process of $P_1$ with the following properties: if this process does not terminate after finitely many steps, then it produces an infinite descending $\Diamond$-tower; if this process terminates after finitely many steps, then it produces a finite tuned refinement of $P_1$ contained in $V$, hence implying that the subalgebra generated by $P_1$ is finite.
    \par \textbf{Step 1}:
    For $\emptyset\neq\mathcal{A}\subseteq P_1$, define
    \[
    \phi_1 (\mathcal{A})= 
    \bigcap_{A \in \mathcal{A}} \Diamond A \ \cap
    \bigcap_{B \in (P_1-\mathcal{A})}\Box (X-B)
    \]
    and let
    \[
    T_1=\Box\emptyset \cup\bigcup_{\emptyset \neq\mathcal{A}\subseteq P_1}(\phi_1(\mathcal{A}) \cap \Box\phi_1(\mathcal{A})).
    \]
    \textbf{Claim.} The following hold:
    \begin{itemize}
    \item [(i)] $T_1 \cap \Diamond(X-T_1)=\emptyset$
    \item [(ii)] $X-T_1 \subseteq \Diamond T_1$.
    \end{itemize}

    \begin{proof}
    \begin{itemize}
    
  \item[(i)] If $x\in T_1$, then either $x \in \Box \emptyset$ or $x \in \phi_1(\mathcal{A}) \cap \Box\phi_1(\mathcal{A})$ for some $\emptyset\neq\mathcal{A}\subseteq P_1$, in which case
  $x \in \Box\Box\phi_1(\mathcal{A})$, and subsequently $x \in \Box T_1$.
    In either case, $x \notin \Diamond(X-T_1)$.
    
  \item[(ii)] Let $x_1 \in X-T_1$. Since $\Box \emptyset \subseteq T_1$, we can assume $R(x_1)\neq \emptyset$.
    Hence, the set $\mathcal{A}_1 = \{A \in P_1 \mid x_1 \in \Diamond A \}$ is nonempty, and $x_1 \in \phi_1(\mathcal{A}_1)$.
    But $x_1 \notin \Box \phi_1(\mathcal{A}_1)$, so there is some $x_2 \in R(x_1)$ such that $x_2 \notin \phi_1(\mathcal{A}_1)$.
    Again, set $\mathcal{A}_2=\{A \in P_1\mid x_2 \in \Diamond A\}$.
    Since $x_2 \in \Diamond A$ implies $x_1 \in \Diamond A$, it follows that
    $\mathcal{A}_2\subsetneq \mathcal{A}_1$ .
    If $x_2 \notin T_1$, then we likewise find some $x_3 \in R(x_2)$ and $\mathcal{A}_3\subsetneq \mathcal{A}_2$ with $x_3 \in \phi_1(\mathcal{A}_3)$. 
    Continue in this manner to get some $x_n \in T_1$ with 
    $x_1Rx_2R\dots Rx_n$, which exists by the well-foundedness of $\mathcal{P}(P_1)$. Thus $x_1 \in \Diamond T_1$, as desired.
    \end{itemize}

    \end{proof}

    Now, if $X-T_1$ is nonempty, we split $P_1$ among $T_1$ and $X-T_1$. Further split $P_1 \restriction T_1$ among 
    \[
    \{\Box\emptyset\}\cup\{\phi_1(\mathcal{A})\mid \emptyset\neq\mathcal{A} \subseteq P_1\}
    \]
    which is a finite partition of $T_1$ contained in $V$. Let $F_1$ denote the result of this splitting.
    \par
    The partition $F_1$ is tuned over $(T_1, R\restriction T_1)$.
        To see why, let $A,B \in F_1$ such that there is an $a \in A$ and $b \in B$ with $aRb$. Then since $a,b \in T_1$, we see that $a,b \in \phi_1(\mathcal{A})$ for some $\mathcal{A} \subseteq P_1$ with $A,B \in \mathcal{A}$.
        Let $c \in A$ and note that $c \in \phi_1(\mathcal{A})$, which follows from the above splitting.
        Hence, $cRd$ for some $d \in B \in \mathcal{A}$ by definition of $\phi_1(\mathcal{A})$, as desired.

    Thus, if $X=T_1$, then we are done and the process ends. If not, let $X_2=X-T_1$. 
    By lemma \ref{prop:exists_tuned_wrt}, we obtain a finite refinement $P_2\subseteq V$ of $P_1\restriction X_2$ which is tuned with respect to $F_1$.
    This concludes step 1.
    \par
    \textbf{Step n+1}: Given that we have obtained $X_{n+1}$ and $P_{n+1}$ from step $n$, we repeat this process with respect to these sets.
    This step will be similar to step 1 in design, but will contain a number of finer differences that warrants an explicit description.
    For $\emptyset\neq\mathcal{A}\subseteq P_{n+1}$, define
    \[
    \phi_{n+1} (\mathcal{A})= 
    \bigcup_1^n T_i\cup\left[ 
    \bigcap_{A \in \mathcal{A}} \Diamond A \ \cap
    \bigcap_{B \in (P_{n+1}-\mathcal{A})}\Box (X-B)
    \right]
    \]
    and let
    \[
    T_{n+1}=X_{n+1}\cap \left[ \Box\bigcup_1^n T_i \cup \bigcup_{\emptyset \neq\mathcal{A}\subseteq P_1}
    \phi_{n+1}(\mathcal{A}) \cap \Box\phi_{n+1}(\mathcal{A}) \right].
    \]
    \textbf{Claim.} 
    \begin{itemize}
    \item[(i)] $T_{n+1} \cap \Diamond(X_{n+1}-T_{n+1})=\emptyset$
    \item[(ii)] $X_{n+1}-T_{n+1} \subseteq \Diamond T_{n+1}$.
    \end{itemize}
    
    \begin{proof}
    \begin{itemize}
    \item[(i)] Let $x \in T_{n+1}$.
        If $x\in X_{n+1}\cap\Box\bigcup_1^n T_i$, then it witnesses nothing in $X_{n+1}$.
        If $x\in X_{n+1} \cap \phi_{n+1}(\mathcal{A}) \cap \Box\phi_{n+1}(\mathcal{A})$, then $x \in \Box\Box\phi_{n+1}(\mathcal{A})$, and it follows that $x \in \Box T_{n+1}$. In any case, $x \notin \Diamond(X_{n+1} - T_{n+1})$.
    \item[(ii)] Let $x_1 \in X_{n+1}-T_{n+1}$. 
    Since $(X_{n+1}\cap\Box\bigcup_1^nT_i) \subseteq T_{n+1}$, we can assume $R(x_1) \cap T_{n+1}$ is nonempty. 
    We omit the rest of the proof, as it is very similar to that of (ii) in the claim in step 1.
    \end{itemize}
    \end{proof}
    \par
    Now, in the case where $T_{n+1}=X_{n+1}$, we can split $P_{n+1}$ among 
    \[
    \mathcal{F}=\{\Box\bigcup_1^nT_i\}\cup \{\phi_{n+1}(\mathcal{A}) \mid \emptyset \neq \mathcal{A} \subseteq P_{n+1}\}
    \]
    which is a finite family of pairwise disjoint members of $V$. Call the result of this splitting $F_{n+1}$. 
    By the same reasoning used to prove that $F_1$ was tuned, $F_{n+1}$ is a tuned partition of $(T_{n+1},R\restriction T_{n+1})$ contained in $V$.
    \par
    Hence, each $F_j$ is a finite tuned partition of $T_j$, and as previously obtained, $F_j$ is tuned with respect to $F_i$ for each $i<j$.
    Moreover, for $i<j$, we see that $A \in F_j$ and $B \in F_i$ implies $\Diamond A \cap B = \emptyset$, so $F_i$ is tuned with respect to $F_j$ as well.
    Consequently, $\bigcup_1^{n+1}F_i$ is a finite tuned refinement of $P_1$ contained in $V$, so the subalgebra generated by $P_1$ is finite. In this case, the process terminates.
    \par
    On the other hand, if $T_{n+1}\neq X_{n+1}$, then let $X_{n+2}=X_{n+1}-T_{n+1}$, and split $P_{n+1}$ among $\{T_{n+1},X_{n+2}\}$.
    Similarly, split $P_{n+1} \restriction T_{n+1}$ among $\mathcal{F}$ to get $F_{n+1}$, a finite tuned refinement of $P_{n+1} \restriction T_{n+1}$ which is contained in $V$.
    Then by lemma \ref{prop:exists_tuned_wrt}, we obtain a finite refinement $P_{n+2}\subseteq V$ of $P_{n+1}\restriction X_{n+2}$ which is tuned with respect to $F_{n+1}$. 
    This concludes step n+1.
    \par
    If this process does not terminate, then the claims proven in each step show that $\{T_1,T_2,\dots\}\subseteq V$ is an infinite descending $\Diamond$-tower, contrary to our assumption. Thus, this process terminates after finitely many steps, implying that it has a finite tuned refinement contained in $V$.
    Since $P_1$ was an arbitrary finite partition of $X$ contained in $V$, it follows that $(X,R,V)$ is tunable and $(V,\Diamond)$ is locally finite.    
\end{proof}

\end{document}
